\section{Related Work}
\label{sec:related_work}

Retinal vessel segmentation has long been a benchmark of medical image analysis~\cite{Staal2004, moccia2018blood}. Deep segmentation networks build on fully convolutional networks~\cite{Long_2015_CVPR} and the U-Net~\cite{ronneberger2015u}, whose skip connections carry high-resolution features to the decoder, and U-Net variants became the primary tool of medical image segmentation~\cite{siddique2021u}. \convnext{}~\cite{liu2022convnet} modernized the CNN with large depthwise kernels and inverted bottlenecks. In the opposite direction, Galdran et al.~\cite{galdran2020littlewnetcouldstateoftheart} showed that a minimal cascaded U-Net, Little W-Net, with about 68,000 parameters performs on par with or better than much larger models on retinal benchmarks when it is carefully trained and evaluated.

The ViT~\cite{dosovitskiy2020image} encodes an image as a sequence of patch tokens processed by self-attention~\cite{vaswani2017attention}. In medical segmentation, Transformers were first combined with CNNs, as in TransUNet~\cite{chen2021transunet}, or used in hierarchical forms such as Swin-Unet~\cite{cao2022swin}. A systematic review of 36 studies that compared ViTs with CNNs in medical image analysis found the transformer-based models often ahead and pre-training important for them~\cite{takahashi2024comparison}; the review was not specific to vessels. The Segment Anything Model (SAM)~\cite{Kirillov_2023_ICCV} introduced a promptable foundation model for segmentation, whose direct use on medical images is limited~\cite{zhang2024segmentmodelmedicalimage}. \medsam{}~\cite{ma2024segment} adapted SAM's ViT-B to medical images with over one million image--mask pairs. Its authors report that it has difficulty with vessel-like branching structures, which they attribute to the ambiguity of the box prompt~\cite{ma2024segment}. \litemedsam{}, the baseline of the CVPR~2024 challenge ``Segment Anything in Medical Images on Laptop'', distilled the encoder of \medsam{} into a TinyViT~\cite{wu2022tinyvit}, a lightweight hybrid of convolution and Transformer layers~\cite{ma2024efficient}.

A plain ViT such as \medsam{} projects each $16 \times 16$ patch to one token, so a vessel one or two pixels wide occupies a small fraction of its token. Chen et al.~\cite{chen2021transunet} attributed coarse segmentations to this loss of low-level detail, which motivated hybrid designs. Patch tokenization is therefore a strong reason to expect plain ViTs to miss thin vessels, although not the only one: plain ViTs also lack the locality bias of convolutions and need large training sets~\cite{dosovitskiy2020image}. We test this expectation on the thin vessels themselves, with a plain ViT evaluated at the scale of its training and compared with CNNs under a training protocol that aims at the same conditions.
